\cotefile{157-2}
\scanpage{227}
\mspage{173}
% archive sheet number: 226
quel est l'axiome \unsure{mettant} en relation \unsure{le|les} \unsure{dérivateur|dérivateurs}
$\mathbb{D}$ avec $W_\infty$, ou avec quelque autre localiseur
fondamental.

Je vais faire à présent un résumé de la
théorie de $\mathrm{HOT}_W$, revue dans l'optique que je
viens de dire : \unsure{On} suppose pour $W$ \uline{les conditions}
$W(1,2,3)$ \uline{sans plus}, sauf mention explicite.

\textcircled{1} \uline{Localiseurs principaux} sur $\mathrm{Cat}_{/S}$, et sur les
\deleted{\unsure{autres}}\inserted{principales} catégories \deleted{\unsure{des diagrammes}} de $S$-catégories (cf diagramme p. \unsure{76}). Ils seront tous induits par les trois
localiseurs familiers
\diagram{\begin{tikzcd}[ampersand replacement=\&,column sep=small,row sep=small]
 W_S^d \& \& W_S^g\\
 \& W_S \arrow[ul,phantom,"\subset"{description,rotate=-45}] \arrow[ur,phantom,"\subset"{description,rotate=-135}] \&
\end{tikzcd}}
sur $\mathrm{Cat}_{/S}$. \deleted{Mais il} Il y a lieu aussi de
regarder $W_S^f$, mais jamais sur $\mathrm{Cat}_{/S}$ tout entier,
mais sur d'autres catégories \unsure{du diagramme|de diagrammes}, et
toujours dans des cas où $W_S^f$ coïncide soit avec
$W_S^d$, soit avec $W_S^g$. Mais le point de vue $W_S^f$ est
souvent plus \inserted{directement} commode, et a l'avantage de
plus d'être symétrique.

\textcircled{a} \deleted{Dans} \uline{$\mathrm{Liss}$}$_W\,S$, on a
\diagram{\begin{tikzcd}[ampersand replacement=\&,row sep=small]
 W_S^f \arrow[d,Rightarrow] \& {}=W_S^d=W_S^{\mathrm{univ}}\\
 W_S^g \arrow[d,Rightarrow] \&\\
 \underline{\hspace{1em}W_S\hspace{1em}} \& .
\end{tikzcd}}

Dualement \deleted{\unsure{dans}} \textcircled{b} \deleted{\unsure{dans}} \uline{$\mathrm{Prop}$}$_W\,S$, on a
\diagram{\begin{tikzcd}[ampersand replacement=\&,row sep=small]
 W_S^f \arrow[d,Rightarrow] \& {}=W_S^g=W_S^{\mathrm{univ}}\\
 W_S^d \arrow[d,Rightarrow] \&\\
 W_S \& {,}
\end{tikzcd}}
\scanpage{228}
\mspage{174}
% archive sheet number: 227
enfin

\textcircled{c} \deleted{dans} \uline{$\mathrm{Fib}$}$_W\,S$, on a
\diagram{\begin{tikzcd}[ampersand replacement=\&,row sep=small,column sep=small]
W_S^f \arrow[r,Leftrightarrow] \arrow[d,Rightarrow] \& W_S^d \arrow[r,Leftrightarrow] \& W_S^g \arrow[r,Leftrightarrow] \& W_S^{\mathrm{univ}}\\
W_S \& {.} \& \&
\end{tikzcd}}

Cette implication $W_S^f\Rightarrow W_S$ est stricte si on ne fait pas
sur $W$ l'axiome $W(4\mathrm{b})$, \unsure{dont} l'utilité ici est
exactement d'\unsure{assurer}
\[
\left|\begin{aligned}
&W_S^f\Longleftrightarrow W_S\quad\text{dans }\underline{\mathrm{Fib}}_W\,S
  \quad\underline{\text{\uuline{si}}}\quad W(4\mathrm{b})\ \text{satisfaite}\\
&\text{et alors les cinq localiseurs envisagés ci-dessus}\\
&\qquad\text{dans }\underline{\mathrm{Fib}}_W\,S\text{ sont tous égaux.}
\end{aligned}\right.
\]

\deleted{\unsure{Mais}} Les relations \deleted{\unsure{entre}} localiseurs \deleted{\unsure{ci-dessus}} dans
$\underline{\mathrm{Liss}}_W\,S$, $\underline{\mathrm{Prop}}_W\,S$, $\underline{\mathrm{Fib}}_W\,S$ se transportent automatiquement aux sous-catégories de celles-ci, plus gén. à \unsure{toutes|toute} \unsure{celles|celle} qui s'envoient par un foncteur
d'\unsure{inclusion|induction} dans une de celles-ci. Dans le diagramme
p. \unsure{75|76}, ttes les catégories, à la seule exception
de $\mathrm{Cat}_{/S}$ tout entier, sont \deleted{dans} en ce cas.

\deleted{\unsure{\textcircled{d}}\illegible{1}}\inserted{Pour} les résultats principaux développés jusqu'à
présent sur $\mathrm{HOT}_W$, quand l'axiome $W(4\mathrm{b})$ est
utilisé, la forme faible \deleted{\unsure{(}}$W(4\mathrm{b}')$ suffit. \deleted{Celle-ci \unsure{assure}} On a
\[
\left|\begin{aligned}
&W_S^f\Longleftrightarrow W_S\quad\text{dans }\underline{\mathrm{Parf}}_W\,S
 \quad\underline{\text{\uuline{si}}}\quad W(4\mathrm{b}')\ \text{satisfaite}\\
&\text{et \inserted{alors} les cinq localiseurs }W_S^f,W_S^d,W_S^g,W_S^{\mathrm{univ}},W_S\\
&\text{sont égaux dans }\mathrm{Parf}_W\,S .
\end{aligned}\right.
\]

\textcircled{2} \uline{Principales catégories de catégories sur} $S$.

Ce sont, à des modifications d'éclairage près, celles
du diagramme p. \unsure{76}. Les catégories $\underline{\mathrm{Fibsc}}\,S$, $\underline{\mathrm{Fib}}\,S$,
$\underline{\mathrm{Cofibsc}}\,S$, $\underline{\mathrm{Cofib}}\,S$ et leurs variantes d'indice $0$,
sont inchangées -- leur rôle est d'être aussi petites
\scanpage{229}
\mspage{175}
% archive sheet number: 228
que \unsure{possibles|possible} pour remplir leur fonction (de description de $\mathrm{HOT}_W(S)$, $\mathrm{HOT}_W^{\unsure{0|o}}(S)$, $\mathrm{HOT}_W^{\mathrm{lc}}(S)$, ou $\mathrm{HOT}_W^{\mathrm{lc}!}(S)$
ci-dessous). La catégorie $\underline{\mathrm{Parf}}\,S$ du diagramme
remplissait une fonction du même type, c'est
pourquoi nous la remplaçons à présent (en \unsure{surenchérissant|sur-enchérissant}) par $\underline{\mathrm{Parf}}_W\,S$. Les autres catégories du
diagramme sont, en plus de $\mathrm{Cat}_{/S}$, $\underline{\mathrm{Liss}}\,S$ et
$\underline{\mathrm{Prop}}\,S$ et leurs variantes à indices $0$, ainsi que
$\underline{\mathrm{Fibhot}}\,S$. (Toutes ces catégories, ainsi que $\underline{\mathrm{Parf}}\,S$, étaient
donc construites par référence à $W_\infty$, non pas
à $W$.) Leur rôle \deleted{était} est d'être aussi grandes
que possible pour donner lieu à certaines
propriétés et relations. Or il se trouve qu'ici la
référence à $W_\infty$ était artificielle, et il y a
lieu de mettre plutôt les catégories
$\underline{\mathrm{Liss}}_W\,S$, $\underline{\mathrm{Prop}}_W\,S$, $\underline{\mathrm{Fib}}_W\,S$, \unsure{1)} en plus des variantes
$\underline{\mathrm{Liss}}_{0W}\,S$, $\underline{\mathrm{Prop}}_{0W}\,S$. Ainsi le diagramme devient :

\margin{NB La nouvelle notation serait plutôt $\mathrm{Fibf}_W\,S$, car il s'agit de $W$-fibrations qu'à présent \unsure{j'appelle|j'appelles} \textquotedblleft faibles\textquotedblright\ (et définies seulement dans $\mathrm{Cat}$, à l'opposé des $W$-fibrations, définissables dans tte catégorie de modèles).}
\diagram{\scriptsize\begin{tikzcd}[ampersand replacement=\&,column sep=.2em,row sep=.55em]
 \& \& \& \underline{\mathrm{Fib}}_W S \arrow[dd,hook] \arrow[ddddlll,bend right=17,gray,"\Phi_0"'] \arrow[ddddrrr,bend left=17,gray,"\Psi_0"] \& \& \&\\
 \& \& \underline{\mathrm{Liss}}_{0W}S \arrow[ur,hook] \arrow[ddd,hook] \& \& \underline{\mathrm{Prop}}_{0W}S \arrow[ul,hook] \arrow[ddd,hook] \& \&\\
 \& \& \& \mathrm{Cat}_{/S} \arrow[d,gray,"C"'] \arrow[ddddlll,bend right=8,gray,"\Phi"'] \arrow[ddddrrr,bend left=8,gray,"\Psi"] \& \& \&\\
 \& \underline{\mathrm{Fib}}_{0W}S \arrow[uur,hook] \arrow[dd,hook] \& \& \underline{\mathrm{Parf}}_W S \arrow[uul,hook] \arrow[uur,hook] \& \& \underline{\mathrm{Cofib}}_{0W}S \arrow[uul,hook] \arrow[dd,hook] \&\\
 \underline{\mathrm{Fibsc}}_{0W}S \arrow[ur] \arrow[dd,hook] \& \& \underline{\mathrm{Liss}}_W S \arrow[uur,hook] \& \deleted{\unsure{\mathrm{Parf}\,S}} \& \underline{\mathrm{Prop}}_W S \arrow[uul,hook] \& \& \underline{\mathrm{Cofibsc}}_{0W}S \arrow[ul] \arrow[dd,hook]\\
 \& \underline{\mathrm{Fib}}\,S \arrow[ur,hook] \& \& \deleted{\unsure{\underline{\mathrm{Parf}}_W S}} \arrow[uu,hook] \& \& \underline{\mathrm{Cofib}}\,S \arrow[ul,hook] \&\\
 \mathrm{Fibsc}\,S \arrow[ur] \& \& \& \& \& \& \underline{\mathrm{Cofibsc}}\,S \arrow[ul]
\end{tikzcd}}
(tous les carrés cartésiens ; les foncteurs notés $\hookrightarrow$ sont pl. fid.)
\scanpage{230}
\mspage{176}
% archive sheet number: 229
où j'ai \unsure{même|aussi} marqué en vert les cinq ``\uline{foncteurs} d'\uline{\unsure{inversion}}''
\[
\Phi,\Psi,\Phi_0,\Psi_0,C,
\]
qui fournissent des quasi-inverses des foncteurs sur
les \inserted{catégories} localisées \unsure{déduits|induits} par les \uline{foncteurs d'inclusion}.

\textcircled{3} \uline{Les quatre types de catégories} $\mathrm{HOT}_W$

Sur chacune des catégories de l'étage, il y a
\uline{deux} localiseurs à regarder
\diagram{\begin{tikzcd}[ampersand replacement=\&,row sep=small,column sep=0pt]
 W_S^f \arrow[d,Rightarrow] \& {}=W_S^g=W_S^d=W_S^{\mathrm{univ}}\\
 W_S \&
\end{tikzcd}}
donnant lieu à deux localisés distincts, notés
\[
\mathrm{HOT}_W^{\mathrm{lc}}(S)\qquad\text{et}\qquad\mathrm{HOT}_W^{\mathrm{lc}!}(S),
\]
il se trouve que (grâce à $\Phi_0$ et $\Psi_0$, \deleted{et $C$ pour $\underline{\mathrm{Parf}}_W$}) ces localisés
``sont les \deleted{\unsure{m}}\unsure{même|mêmes}'' pour les 8 catégories de l'étage,
\deleted{\unsure{y inclus}}\inserted{à l'exception de} $\underline{\mathrm{Parf}}_W\,S$ ; plus exactement les foncteurs
d'inclusion \inserted{(à l'exception peut-être de ceux ayant pour source $\underline{\mathrm{Parf}}_W\,S$)}
induisent des équivalences sur les
deux types de localisés, les équivalences quasi-inverses étant obtenues par $\Phi_0$ (aile gauche) ou
$\Psi_0$ (aile droite). \inserted{Quand on localise \uline{pour} $W_S$ (à l'exclusion de $W_S^f$) on trouve aussi \deleted{des} quasi-inverses resp.}
\deleted{\unsure{ou}} par $C\mid\underline{\mathrm{Liss}}_{0W}\,S$
$C\mid\underline{\mathrm{Prop}}_{0W}\,S$, pour les deux inclusions
\diagram{\begin{tikzcd}[ampersand replacement=\&]
\underline{\mathrm{Liss}}_{\unsure{0W|W}}\,S \& \underline{\mathrm{Parf}}_W\,S \arrow[l,hook] \arrow[r,hook,no head] \& \underline{\mathrm{Prop}}_{0W}\,S .\quad {*}
\end{tikzcd}}
% slanted
\margin{cf plus bas ce qu'on peut dire pour les localisés pour $W_S^f$}
Par raison de symétrie, le plus joli semble de
définir
\[
\left\{\begin{aligned}
\mathrm{HOT}_W^{\mathrm{lc}}(S)&=(\underline{\mathrm{Fib}}_W\,S)(W_S^f)^{-1}\\
\mathrm{HOT}_W^{\mathrm{lc}!}(S)&=(\underline{\mathrm{Fib}}_W\,S)(W_S)^{-1},
\end{aligned}\right.
\]
et on a un foncteur de localisation
\[
\mathrm{HOT}_W^{\mathrm{lc}}(S)\xrightarrow{\unsure{\lambda|{?}}}\mathrm{HOT}_W^{\mathrm{lc}!}(S)
\]
induit par l'inclusion
\[
W_S^f\subset W_S .
\]
\scanpage{231}
\mspage{177}
% archive sheet number: 230
Pour toute sous-catégorie pleine
\[
\unsure{\mathcal{F}|F}\subset\underline{\mathrm{Fib}}_W\,S
\]
qui contient l'une des \deleted{deux}\inserted{trois} sous-catégories
\[
\underline{\mathrm{Fib}}_{0W}\,S,\quad \deleted{\underline{\mathrm{Parf}}_W\,S,}\inserted{\underline{\mathrm{Parf}}_{\unsure{W|0W}}\,S,}
\quad\underline{\mathrm{Prop}}_{0W}\,S,
\]
l'inclusion de $\mathcal{F}$ dans $\underline{\mathrm{Fib}}_W\,S$ définit une équivalence
\[
\mathcal{F}(W_S^{\deleted{\unsure{f}}})^{-1}\approx\mathrm{HOT}_W^{\mathrm{lc}!}(S),
\]
\inserted{et si $\mathcal{F}$ contient l'une des deux sous-catégories $\underline{\mathrm{Fib}}_{0W}\,S$, $\underline{\mathrm{Prop}}_{0W}\,S$ \inserted{(à l'excl. de $\underline{\mathrm{Parf}}_W\,S$)} on trouve aussi}
\[
\mathcal{F}(W_S^f)^{-1}\approx\mathrm{HOT}_W^{\mathrm{lc}\deleted{\unsure{!}}}(S).
\]

Le rez de chaussée contribue deux autres variantes
de $\mathrm{HOT}$, en localisant soit par $W_S^d$ (aile gauche),
soit \deleted{\unsure{par}} $W_S^g$ (aile droite). Les foncteurs d'inclusion
définissent des équivalences pour les localisés soit
pour $W_S^d$ (a. g.), soit pour $W_S^g$ (a. dr.). Le
plus joli (symétrie) semble de définir ces deux
localisés par $\mathrm{Cat}_{/S}$
\[
\mathrm{HOT}_W(S)=(\mathrm{Cat}_{/S})(W_S^d)^{-1},\quad
\mathrm{HOT}_W^{\unsure{0|o}}(S)=(\mathrm{Cat}_{/S})(W_S^g)^{-1}.
\]
Les quasi-inverses pour les équivalences de catégories
localisées \deleted{\unsure{d'in}}déduites par les foncteurs d'inclusion,
sont donnés soit par $\Phi$ (aile gauche), soit par $\Psi$
(aile droite). \deleted{\unsure{M}}

Mais $\Phi$ et $\Psi$ \deleted{\unsure{peuvent se décrire par}}\inserted{fournissent des} quasi-inverses également quand on localise pour $W_S$,
\deleted{\unsure{disons équi}} $W$-équivalences grossières au lieu de le faire
-- pareil qu'à l'étage. On trouve donc au rez
de chaussée, un \inserted{(unique)} troisième localisé, qu'on peut
décrire \deleted{\unsure{symétr}}\deleted{\unsure{le Cat}} plus joliment par
\[
(\mathrm{Cat}_{/S})(W_S)^{-1},
\]
mais il est équivalent à $(\underline{\mathrm{Fib}}_W\,S)W_S^{-1}$ par le
\scanpage{232}
\deleted{\unsure{176}}\mspage{178}
% archive sheet number: 231
foncteur induit par l'inclusion. \inserted{(\unsure{Un|son} qu. inv. étant fourni par $C$.)}
\deleted{On trouve ainsi}\inserted{C'est donc}
$\mathrm{HOT}_W^{\mathrm{lc}!}(S)$. On trouve ainsi, par le diagramme
\inserted{\unsure{comm}} de catégories avec localiseur
\diagram{\begin{tikzcd}[ampersand replacement=\&,row sep=small,column sep=small]
 \&(\underline{\mathrm{Fib}}_W\,S,W_S^f)\arrow[dl]\arrow[dr]\arrow[dd]\&\\
 (\mathrm{Cat}_{/S};W_S^d)\arrow[ddr]\&\&(\mathrm{Cat}_{/S},W_S^g)\arrow[ddl]\\
 \&(\underline{\mathrm{Fib}}_W\,S,W_S)\arrow[d,thick]\&\\
 \&(\mathrm{Cat}_{/S},W_S)\&\deleted{\unsure{S}}
\end{tikzcd}}
(où la flèche verticale courte induit une équivalence sur les localisées), le diagramme \unsure{comm.}
entre les catégories $\mathrm{HOT}_W$ liées à $S$ :
\diagram{\begin{tikzcd}[ampersand replacement=\&,column sep=small]
 \& \mathrm{HOT}_W^{\mathrm{lc}}(S)\arrow[dl,"i"']\arrow[dd,"\lambda"']\arrow[dr,"j"] \&\\
 (*)\quad \mathrm{HOT}_W(S)\arrow[dr,"\rho"'] \& \& \mathrm{HOT}_W^{\unsure{0|o}}(S)\;(\simeq\mathrm{HOT}_W(S^{\unsure{o|0}}))\arrow[dl,"\sigma"]\\
 \& \mathrm{HOT}_W^{\mathrm{lc}!}(S)\& {,}
\end{tikzcd}}
où $\rho,\sigma,\lambda$ sont les foncteurs de localisation,
et où $i,j$ sont ressentis comme des
foncteurs d'inclusion. Ce diagramme $(*)$
\inserted{semble} épuiser les foncteurs naturels qu'on \deleted{\unsure{peut}}\inserted{\unsure{relient|relier}}
\deleted{établir} entre elles les quatre catégories envisagées.
\scanpage{233}
\mspage{179}
% archive sheet number: 232
Les propriétés des foncteurs $\Phi,\Psi,C$ qui ont été
utilisées, pour \deleted{\unsure{obtenir les}}\inserted{fournir} tous les quasi-inverses
voulus des foncteurs (entre cat. localisées) induits par les foncteurs d'inclusion du diagramme, sont les suivantes.

\textcircled{a} On a \unsure{un} lien fonctoriel dans $\mathrm{Cat}_{/S}$
\[
X\xrightarrow{i_X}\Phi(X),\qquad \underline{\underline{\unsure{j|i}_X\in W_S^{\unsure{d|1}}}} .
\]
\deleted{\unsure{qui est dans} $W_S^d$,} Quand $X$ est dans $\underline{\mathrm{Fib}}_W\,S$,
$\Phi(X)$ est dans $\underline{\mathrm{Fibsc}}_0\,\deleted{\unsure{S}}S\deleted{\unsure{S}}$,
\deleted{\unsure{(sa restriction), donc}}\inserted{donc $i_X\in W_S^f$ ($=W_S^d$ dans $\underline{\mathrm{Fib}}_W\,S$)}.
\inserted{Ainsi} $\Phi$ induit le foncteur $\Phi_0$ du diagramme.

Dualement, on a
\[
X\mathrel{\overset{i_X}{\unsure{\longrightarrow|\text{---}}}}\Psi(X),
\qquad\underline{i_X\in W_S^g},
\]
et quand $X$ est dans $\underline{\mathrm{Fib}}_W\,S$, $\Psi(X)$ est dans
$\underline{\mathrm{Cofibsc}}_0\,S$, d'où le foncteur $\Psi_0$.

\textcircled{b} On a un lien fonctoriel
\[
\left|\begin{aligned}
&X\mathrel{\overset{i_X}{\unsure{\text{---}|\longrightarrow}}}C(X),\qquad\underline{i_X\in W_S}.\\
&\text{et }C(X)\in\mathrm{Ob}\,\underline{\mathrm{Parf}}_W\,S .
\end{aligned}\right.
\]
\deleted{De plus, quand $X$ est dans $\underline{\mathrm{Fib}}_W\,S$, on a
\[
i_X\in W_S^f\quad\text{si}\quad X\in\mathrm{Ob}\,\underline{\mathrm{Fib}}_W\,S .
\]}
Mais on fera attention qu'en l'absence de l'axiome
\margin{$\vert\vert\vert$}$W(4\mathrm{b})$ sur $W$, on ne peut \uline{pas} affirmer $i_X\in W_S^f$,
même si $X$ est dans $\underline{\mathrm{Fib}}_W\,S$\deleted{\unsure{.}} (voire, dans $\underline{\mathrm{Parf}}_W\,S$),
seul cas (de toutes façons) où on est en droit de s'y
attendre. Je note à ce propos le diagramme \uline{cartésien},
pour $f:X\mathrel{\unsure{\text{---}|\longrightarrow}}\deleted{\unsure{C}}Y$ dans $\mathrm{Cat}_{/S}$
\diagram{\begin{tikzcd}[ampersand replacement=\&,column sep=large]
C(Y)_{\unsure{s|S}}\arrow[d,"\pi_Y^s"'] \& C(X)_s\arrow[l,"C(f)_s"']\arrow[d,"\pi_X^s"]\\
Y \& X\arrow[l,"f"']\;{,}
\end{tikzcd}}
\margin{flèches verticales dans $\mathrm{Parf}_W\subset\mathrm{Fib}_W$}
qui montre par contre que \deleted{si $i_X,i_Y\in W_S^f$, \unsure{alors} $f\in W_S$,}
\scanpage{234}
\mspage{180}
% archive sheet number: 233
\deleted{pour \unsure{le|un} \unsure{S-morphisme|S-équivalence} $f:X\longrightarrow Y$ qui est dans $W$, alors}
l'on ne peut pas s'attendre que $f\in W_S\Rightarrow C(f)\in\underline{\underline{W_S^f}}$
sans disposer de l'axiome $W(4\mathrm{a}')$ (ou d'un
axiome de force comparable \deleted{\unsure{)}}, tel $W(\deleted{\unsure{b}}4\mathrm{b})$).

Ce sont ces faits \deleted{qui} négatifs (en l'absence
de $W(4\mathrm{b})$) qui sont cause que je ne peux rien
dire de particulier sur les foncteurs induits
par \uline{localisation fine} $(W_S^f)^{-1}$ des deux foncteurs
d'inclusion
\diagram{\begin{tikzcd}[ampersand replacement=\&,column sep=small]
\underline{\mathrm{Liss}}_0\,S \& \underline{\mathrm{Parf}}_{\unsure{\infty|W}}\,S \arrow[l]\arrow[r]\arrow[dl] \& \underline{\mathrm{Prop}}_0\,S,\\
{} \& \&
\end{tikzcd}}
ou ce qui revient au \unsure{$=$|même}, de celui induit par le
composé commun de ceux-là avec les inclusions
dans $\underline{\mathrm{Fib}}_W\,S$
\diagram{\begin{tikzcd}[ampersand replacement=\&,row sep=small,column sep=small]
 \&\underline{\mathrm{Parf}}_{\unsure{w|W}}\,S\arrow[dl]\arrow[dr]\arrow[dd]\&\\
 \underline{\mathrm{Liss}}_0\,S\arrow[dr]\&\&\underline{\mathrm{Prop}}_{\unsure{1|0}}\,S\arrow[dl]\\
 \&\underline{\mathrm{Fib}}_W\,S\&{;}
\end{tikzcd}}
en d'autres termes, je ne peux rien dire sur
les propriétés de fidélité, \deleted{\unsure{voire}} du foncteur
\[
(\underline{\mathrm{Parf}}_W\,S)(W_S^f)^{-1}\mathrel{\unsure{\longrightarrow|\text{---}}}
(\underline{\mathrm{Fib}}_W\,S)(W_S^f)^{-1},
\]
et encore moins que ce soit une \unsure{équiv.} de
catégories. Donc, en l'absence de $W(4\mathrm{b})$,
pour les sous-catégories pleines \inserted{$\unsure{\mathcal{F}|F}$} de $\underline{\mathrm{Fib}}_W\,S$ qui
contiennent $\underline{\mathrm{Parf}}_W\,S$, sans contenir ni $\underline{\mathrm{Liss}}_0\,S$ ni
$\underline{\mathrm{Prop}}_0\,S$, il se pourrait a priori que les
localisées
\[
\mathcal{F}(W_S^f)^{-1}
\]
soient toutes différentes. Je pense notamment
à $\underline{\mathrm{Parf}}_W\,S$, aussi \unsure{intéressante|intéressant} en elle-même, après
tout, que $\underline{\mathrm{Parf}}_{\unsure{w|\infty|W}}\,S$. Donc en plus des quatre
localisées du diagramme $(*)$ p. 178, il y a toute
une série infinie ``tératologique'' de localisées
\scanpage{235}
\mspage{181}
% archive sheet number: 234
qui refusent de se laisser ``\unsure{toucher|boucher}'' par $\Phi_0$ ou $\Psi_0$, et
pour lesquelles \unsure{même|mêmes} $C$ n'est d'aucun secours,
faute de l'axiome $W(4\mathrm{b})$.

\textcircled{4} \uline{Cas où l'axiome} $W(4\mathrm{b})$ \uline{est satisfait}

Alors les perplexités que je viens de dire disparaissent. On a aussitôt, en vertu de $W(4\mathrm{b})$ :
\[
\left|\!\left|\!\left|\begin{aligned}
&\text{Si }X\in\mathrm{Ob}\,\underline{\mathrm{Fib}}_W\,S,\quad\text{alors }i_X:X\longrightarrow C(X)\quad\text{est dans }W_S^f\\
&\qquad\text{(\uline{en supposant l'axiome $W(4\mathrm{b})$})}
\end{aligned}\right.\right.\right.
\]
et $C$ fournit des quasi-inverses pour les foncteurs
``d'inclusion'' entre localisés fins (par $W_S^f$) envisagés à la
page précédente. \deleted{Il en fournit} un également
pour le foncteur $\lambda$ envisagé plus haut ($(*)$ page 178)
\diagram{\begin{tikzcd}[ampersand replacement=\&,row sep=.15em]
(\underline{\mathrm{Fib}}_W\,S)(W_S^{\deleted{\unsure{f}}})^{-1}\arrow[r]\arrow[d,phantom,"\Vert\ {\mathrm{def}}"] \&
(\mathrm{Cat}_{/S})(W_S^{-1})\arrow[d,phantom,"\deleted{\unsure{\Vert}}\,\unsure{\wr\wr|{??}}"]\\
\mathrm{HOT}_W^{\mathrm{lc}}(S) \& \mathrm{HOT}_W^{\mathrm{lc}!}(S)\;{,}
\end{tikzcd}}
lequel est donc une équivalence. Une autre façon,
plus évidente, de dire la $=$ chose, c'est que sur
\inserted{chacune} des catégories de l'étage les \inserted{deux} localiseurs $W_S^f$ et $W_S$
sont désormais identiques (puisqu'elles sont toutes
contenues dans $\underline{\mathrm{Fib}}_W\,S$, où elles le sont justement par
l'axiome $W(4\mathrm{b})$), donc le foncteur $\lambda$ entre
ces deux localisés (définition originelle de $\lambda$,
dans le cas $\unsure{\mathcal{F}|F}=\underline{\mathrm{Fib}}_W\,S$, cf \unsure{page} 176) est même
le \uline{foncteur identique}, les deux catégories \inserted{localisées} étant \uline{égales} :
\[
\mathrm{HOT}_W^{\mathrm{lc}}(S)=\mathrm{HOT}_W^{\mathrm{lc}!}(S)\qquad !
\]
Ainsi le diagramme $(*)$ p. 178 prend la forme
\diagram{\begin{tikzcd}[ampersand replacement=\&,column sep=small]
 \&\mathrm{HOT}_W^{\mathrm{lc}}(S)\arrow[dl,"i"]\arrow[dr,"j"'] \&\\
 (*)\quad\mathrm{HOT}_W(S)\arrow[ur,bend left=17,"\rho"] \& \&\mathrm{HOT}_W^{\unsure{0|o}}(S)\arrow[ul,bend right=17,"\sigma"'] .
\end{tikzcd}}
\scanpage{236}
\mspage{182}
% archive sheet number: 235
Quand on interprète non seulement les deux catégories
latérales $\mathrm{HOT}_W(S)$, $\mathrm{HOT}_W^{\unsure{0|o}}(S)$, mais aussi celle du
milieu, comme des localisés de $\mathrm{Cat}_{/S}$, par les
localiseurs $W_S^d,W_S^g$ et $W_S$ respectivement, les foncteurs
$\rho,\sigma$ sont ceux déduits des \inserted{deux} inclusions
\diagram{\begin{tikzcd}[ampersand replacement=\&,column sep=small,row sep=small]
 W_S^d\&\& W_S^g\\
 \&W_S\arrow[ul,phantom,"\subset"{description,rotate=-45}]\arrow[ur,phantom,"\subset"{description,rotate=-135}]\&
 \text{comme parties de }\mathrm{Fl}(\mathrm{Cat}_{/S}),
\end{tikzcd}}
tandis que les foncteurs $i,j$ sont déduits, \deleted{\unsure{par}}
par passage aux localisés, d'un seul et même
foncteur
\[
\mathrm{Cat}_{/S}\xrightarrow{C}\mathrm{Cat}_{/S},
\]
ce foncteur ayant la vertu d'\deleted{\unsure{être}}\inserted{appliquer} $W_S$ dans
$W_S^d\cap W_S^g$.
\[
f\in W_S\underset{\substack{\text{moyennant}\\W(4\mathrm b)}}{\Longrightarrow}
 C(f)\in W_S^d\cap W_S^g\quad(\text{et }{=}\ \in W_S^{\mathrm{univ}})
\]
(\unsure{Et} pour cause, puisque $C$ se factorise par
$\underline{\mathrm{Fib}}_W\,S$ (et $=$ par $\underline{\mathrm{Parf}}_W\,S$), où $W_S$ est
identique à $W_S^d,W_S^g,W_S^{\mathrm{univ}}$. \deleted{\unsure{)}} Que $f\in W_S$ implique $C(f)\in W$ est indépendant de l'axiome
$W(4\mathrm b)$, et se voit sur le diagramme
\diagram{\begin{tikzcd}[ampersand replacement=\&,row sep=small]
X\arrow[r,no head,"f"]\arrow[d,"i_X"']\&Y\arrow[d,"i_Y"]\\
C(X)\arrow[r,no head,"C(f)"']\& C(Y){,}
\end{tikzcd}}
où les flèches $i_X,i_Y$ sont dans $W_S$.)

Ceci dit, il est vrai que les \deleted{\unsure{deux}} couples
\[
(\rho,i),\quad(\sigma,j)
\]
sont des couples de foncteurs adjoints, donc
($\rho$ et $\sigma$ étant des foncteurs de localisation) que $i$ et
$j$ sont \uline{pleinement fidèles}. La formule d'adjonction
peut se lire comme
\[
\boxed{\mathrm{Hom}_{W_S^g}(X,CY)\simeq\mathrm{Hom}_{W_S^d}(X,CY)
 \simeq\mathrm{Hom}_{W_S}(X,Y)}\tag*{$(*)$}
\]
\scanpage{237}
\mspage{183}
% archive sheet number: 236
pour $X,Y$ dans $\mathrm{Cat}_{/S}$ sans plus.
\margin{Un $\mathrm{Hom}_\Sigma(X,Y)$ indique le $\mathrm{Hom}$ dans $(\mathrm{Cat}_{/S})\Sigma^{-1}$.}

\uline{Corollaire} Si $X$ dans $\unsure{\mathrm{Cs}|\mathrm{Cat}}_{/S}$, $Y$ dans $\underline{\mathrm{Fib}}_W\,S$, alors
\diagram{\begin{tikzcd}[ampersand replacement=\&,column sep=tiny]
\mathrm{Hom}_{W_S^g}(X,Y)\deleted{\unsure{\simeq}}\arrow[dr,"\sim"] \& \&\mathrm{Hom}_{W_S^d}(X,Y)\arrow[dl,"\sim"']\\
\&\mathrm{Hom}_{W_S}(X,Y)\&
\end{tikzcd}}
sont \deleted{des isomorphismes}\inserted{bijectifs}.

En effet, ``$Y$ dans $\underline{\mathrm{Fib}}_W(S)$ implique que $Y\xrightarrow{i_Y}C(Y)$ est
dans $\deleted{\unsure{W_S^f=}}W_S^{\mathrm{univ}}\subset W_S^d\cap W_S^g$, donc
\[
\mathrm{Hom}_{W_S^g}(X,Y)\xrightarrow{\sim}\mathrm{Hom}_{W_S^g}(X,CY)
\mathrel{\underset{\substack{\vert\\\text{adjonction}}}{\simeq}}\mathrm{Hom}_{W_S}(X,Y),
\]
et itou pour $\mathrm{Hom}_{W_S^d}(X,Y)$.

\textcircled{5} \deleted{\unsure{Que reste}} \uline{Perplexités sur les variantes} \inserted{\uline{faibles}} \uline{de l'axiome} $W(4\mathrm b)$.

Que reste-t-il de toutes ces belles choses (plus
que trois catégories $\mathrm{HOT}$, au lieu de $4+\infty$ !) quand
on \deleted{\unsure{ne}} suppose plus que $W(4\mathrm b')$ au lieu de $W(4\mathrm b)$ ?
i.e. l'égalité de $W_S$ et $W_S^f$ dans $\underline{\mathrm{Parf}}_W\,S$ ?
% slanted
\margin{\uline{NB} Si $W$ satisfait l'axiome des limites, alors $W(4\mathrm b')$ implique $W(4\mathrm b)$, \uline{cf} \uline{théorème} récapitulatif dans \uline{XVI} 4.4.13 (p. 58--62), notamment $2^\circ$). \unsure{La condition|Les conditions} $W(4\mathrm b)$ est la condition B \unsure{dudit} énoncé, et $W(4\mathrm b')$ la condition D$'$. (Elle est impliquée par $W(4\mathrm a)\equiv C$, ce que je ne \unsure{soupçonnais} nullement en écrivant \uline{XII} \ldots )\unsure{.}}
Finalement, il me semble (contrairement à ce que
j'avais cru voir) qu'il ne reste presque rien.
On ne sait \inserted{alors} toujours pas, pour $X$ dans $\underline{\mathrm{Fib}}_W\,S$,
\[
i_X:X\longrightarrow C(X)
\]
dans $W_S^f$, et non seulement dans $W_S$ ; on
le sait que pour $X\in\underline{\mathrm{Parf}}_W\,S$, mais
\unsure{alors|là}, \unsure{on en a moins} -- \unsure{pas besoin|peu besoin}, car l'utilité
du foncteur $C(X)$ et de l'inclusion $i_X:X\mathrel{\unsure{\text{---}|\longrightarrow}}C(X)$,
c'est ``d'améliorer'' $X$ et de le rendre
parfait, voire même \unsure{un|en|$W$-} parfait. Mais s'il l'est
déjà, c'est plus la peine. Il y aura, comme
\scanpage{238}
\mspage{184}
% archive sheet number: 237
avant, 4 localisées \unsure{éventuellement|essentiellement} distinctes $\mathrm{HOT}_W,\mathrm{HOT}_W^{\unsure{0|o}}$,
$\mathrm{HOT}_W^{\mathrm{lc}},\mathrm{HOT}_W^{\mathrm{lc}!}$, plus une infinité d'autres \deleted{pour les}
$\unsure{\mathcal{F}|F}(W_S^{\unsure{f}})^{-1}$ pour $\mathcal{F}$ entre $\underline{\mathrm{Parf}}_W\,S$ et $\underline{\mathrm{Fib}}_W\,S$ (cf page 180),
avec le seul bénéfice qu'on peut maintenant exclure
le cas $\mathcal{F}=\underline{\mathrm{Parf}}_W\,S$, car dans ce cas $\mathcal{F}(W_S^f)^{-1}=\mathcal{F}(W_S^{\deleted{\unsure{f}}})^{-1}
\approx\mathrm{HOT}_W^{\mathrm{lc}!}(S)$ fait partie des quatre canoniques. La
chose décisive, c'est qu'on n'a toujours aucun moyen
de prouver que
\[
\lambda:\mathrm{HOT}_W^{\mathrm{lc}}(S)\longrightarrow\mathrm{HOT}_W^{\mathrm{lc}!}(S)
\]
est une équivalence, et je doute que \deleted{\unsure{cela le soit}} ce soit
le cas. Il faudrait regarder dans \inserted{$W=W_{\mathbb D}$} \unsure{les|le} $\mathbb{D}$-\unsure{équivalences|équivalence},
où $\mathbb{D}$ est \unsure{le|un} dérivateur $\mathrm{Der}(k\text{-}\mathrm{Mod})$, $k$ un
anneau donné (p.ex $k=\mathbb{Z}$, ou $k$ un corps\inserted{$\mathbb{Z}/p\mathbb{Z}$ ou $\mathbb{Q}$} \ldots).

\deleted{On trouve} Je ne vais pas le faire maintenant
\deleted{\unsure{cependant !}}\deleted{\unsure{HOT}} (Il est vrai que l'axiome $W(4\mathrm b')$ n'est
pas non plus satisfait dans ce cas \ldots)

Si dans l'énoncé de l'axiome, on remplaçait
$\underline{\mathrm{Parf}}_{\unsure{w|W|\infty}}\,S$ par $\underline{\mathrm{Parf}}_W\,S$, rien d'essentiel ne serait
changé -- tout au plus je remplacerais $\underline{\mathrm{Parf}}_w\,S^{\unsure{\prime|\cdot}}$ par
$\underline{\mathrm{Parf}}_W\,S$ dans le brillant diagramme p. 175. Finalement, je suis à présent convaincu qu'il faut
tout ou rien : ou bien \uline{quatre} catégories $\mathrm{HOT}(S)$
pour $W$, en se contentant de $W(1,2,3)$, ou
bien rien que trois, i.e. pas de $\mathrm{HOT}_W^{\mathrm{lc}!}(S)$ --
et pour ça, il ne semble pas qu'on puisse se
contenter de moins de $W(4\mathrm b)$.

\deleted{Que le foncteur de localisation
\[
\lambda:\mathcal{F}(W_S^f)^{-1}\mathrel{\unsure{\text{---}|\longrightarrow}}\mathcal{F}\,W_S^{-1}
\]
soit une équivalence sur les \unsure{6} catégories de l'étage qui}
\scanpage{239}
\mspage{185}
% archive sheet number: 238
Il suffirait, il est vrai, de supposer que $W_S=W_S^f$ sur
l'une quelconque des sept catégories de l'étage autres
que $\underline{\mathrm{Parf}}_W\,S$. Cela impliquerait alors que \deleted{pour} les
six autres, on aura $\overline{W_S}=\overline{W_S^f}$ (mêmes saturés
forts dans $\mathrm{Fl}(\unsure{\mathcal{F}|F})$\unsure{,}). \inserted{Un des deux} \unsure{les} plus faibles \deleted{\unsure{de ces}} axiomes
en question serait :

Si $f:X\longrightarrow Y$ est un foncteur \uline{cartésien} de catégories
\inserted{(et $W$-fibrées$^{*}$)} \uline{fibrées} sur une même $S$, si $f\in W$, alors $f_s\in W$
pour \deleted{tout} $s\in S$,
% slanted
\margin{``$W$-fibrées'' signifiant ici que pour $u\in\mathrm{Fl}\,S$, les foncteurs $u^*$ de chgt de base $X_{s'}\mathrel{\unsure{\text{---}|\longrightarrow}}X_s$, $Y_{s'}\mathrel{\unsure{\text{---}|\longrightarrow}}Y_s$ sont dans $W$}
(et \inserted{l'énoncé} la version duale, pour les foncteurs cocartésiens de
catégories cofibrées). \deleted{\unsure{Je crois}} Sauf erreur, cet axiome
doit déjà impliquer que $\overline W$ satisfait la version
$W(4\mathrm a)$ la plus forte de toutes (a fortiori $W(4\mathrm b)$)
des axiomes $W(4)$. Cet axiome me paraît donc à
présent beaucoup plus raisonnable que $W(4\mathrm b')$, et
ceci d'autant plus qu'il se réfère directement à la
description de $\mathrm{HOT}^{\mathrm{lc}}_W(S)$ en termes de foncteurs
(ou de ``pseudofoncteurs'' si on ne veut supposer les catégories
scindées et les foncteurs compatibles aux scindages)
$S^{\unsure{o|0}}\mathrel{\unsure{\longrightarrow|\text{---}}}\mathrm{Cat}$. D'ailleurs, bien qu'il n'en ait pas
l'air, cet axiome \uline{est} autodual, puisque la
version duale, pour $S$ fixée, est la version directe
pour $S^o$ ! Cela rend cet axiome encore plus
appétissant.

Donc pour formuler $W(4\mathrm a')$, j'ai envie
à présent de laisser tomber $\underline{\mathrm{Parf}}_W$, et de faire pareil.
Si on garde deux versions $W(4\mathrm a')$ (l'ancienne)
et $W(4\mathrm a'')$ (la nouvelle), l'une et l'autre ont
la vertu d'impliquer $W(4\mathrm b)$. Et $W(4\mathrm b')$ (nouvelle
version) doit impliquer $W(4\mathrm a)$ (version ultra-forte)
sinon pour $W$, du moins pour $\overline W$.

Mes idées étaient \unsure{sérieusement|drôlement|sûrement|terriblement} embrouillées
\scanpage{240}
\mspage{186}
% archive sheet number: 239
jusqu'à aujourd'hui même, car j'avais l'idée
que, grâce au mirifique foncteur $C$, la \unsure{catégori|catégorie}
$\underline{\mathrm{Parf}}_W\,S$, comme les sept autres catégories de l'étage,
pouvait servir à définir $\mathrm{HOT}_W^{\mathrm{lc}}(S)$. Or il n'en est
rien, semble-t-il, à moins de \deleted{\unsure{faire}}\inserted{supposer} l'axiome
$W(4\mathrm b)$, ou \deleted{\unsure{bien}} variante $W(4\mathrm b')$ (nouvelle
version). L'ancienne version de $W(4\mathrm b')$ \deleted{\unsure{était}}\inserted{s'avère}
parfaitement inutilisable.

A ce propos, je n'ai tjrs pas répondu clairement
pourtant, par un ``oui'' ou un ``non'', à la question :

\uline{Question} : les inclusions
\diagram{\begin{tikzcd}[ampersand replacement=\&,column sep=small,row sep=small]
 \&\underline{\mathrm{Liss}}_{0W}\,S\\
 \underline{\mathrm{Parf}}_W\,S\arrow[ur]\arrow[dr]\&\\
 \&\underline{\mathrm{Prop}}_{0W}\,S
\end{tikzcd}}
induisent-elles, même en l'absence de $W(4\mathrm b)$
(ou de ses variantes affaiblies) des équivalences
pour les localisés fins (par $W_S^f$) ? Même
question pour $\underline{\mathrm{Parf}}_{\unsure{\infty|W}}\,S$ au lieu de $\underline{\mathrm{Parf}}_W\,S$.* 
% slanted
\margin{OK en présence de l'axiome $W(6\mathrm{bis})$ en ce \uline{\unsure{lieu|sens}}, cf p. 192}
Même question pour
\[
\underline{\mathrm{Bifib}}\,S\subset\underline{\mathrm{Parf}}_W\,S,
\]
la catégorie des catégories \deleted{\unsure{S}} sur $S$ qui sont à
la fois $\mathrm{Cat}$-fibrées et $\mathrm{Cat}$-cofibrées. \deleted{\unsure{Si}}

Je ne me rends pas compte, \deleted{si les} même en
admettant (cette fois) l'axiome $W(4\mathrm a)$, si
$\underline{\mathrm{Bifib}}\,S$ est assez grosse pour pouvoir décrire tous
les types d'homotopie \inserted{relatifs} loc. const. sur $S$, i.e. si
\[
(\underline{\mathrm{Bifib}}\,S)(W_S^f)^{-1}\mathrel{\unsure{\text{---}|\longrightarrow}}
(\underline{\mathrm{Fib}}_W\,S)(W_S^f)^{-1}
\]
est au moins ess. surjectif, sinon une équivalence de
catégories. Il manque un foncteur
\[
\mathrm{Cat}_{/S}\mathrel{\unsure{\text{---}|\longrightarrow}}\underline{\mathrm{Bifib}}\,S
\qquad\text{ou du moins}\qquad
\underset{\inserted{\substack{\text{ou}\\\underline{\unsure{\mathrm{Fis}|\mathrm{Fib}}}_W\,S\,?}}}{\underset{\llcorner\!\rule{2em}{.35pt}\!\lrcorner}{\underline{\mathrm{Fib}}_W\,S}}
\mathrel{\unsure{\text{---}|\longrightarrow}}\underline{\mathrm{Bifib}}\,S,
\]
qui \deleted{\unsure{joue}} puisse jouer le rôle de $\Phi$ et de $\Psi$ \uline{simultané}%
\scanpage{241}%
\mspage{187}%
% archive sheet number: 240
\uline{ment}. Je n'ai pas même vérifié que $\Phi$ ne transforme \uline{pas} \deleted{\unsure{non plus}} objets de $\underline{\mathrm{Prop}}_W\,S$ ou de $\underline{\mathrm{Prop}}_{0W}\,S$ en
objets $W$-propres (donc $W$-parfaits, donc \unsure{de toutes façons|de la même façon}
$W$-lisses et $=$ $W_{\unsure{\infty|w}}$-lisses) sur $S$.*
% slanted
\margin{Voir réponse à cette intéressante question page 192.}
\deleted{\unsure{Mettons|Admis} que ce soit \unsure{faux}} \inserted{Si c'était pourtant vrai} sur $\underline{\mathrm{Prop}}_W\,S$, \deleted{\unsure{et de fortes raisons induisent} \unsure{à que le}} \inserted{ça signifierait que l'image essentielle}
de $\underline{\mathrm{Prop}}_W\,S$ dans $\mathrm{HOT}_W(S)$
$\bigl(=(\mathrm{Cat}_{/S})(W_S^d)^{-1}\bigr)$ est contenue dans $\mathrm{HOT}_W^{\mathrm{lc}}(S)$.
Prenant $S=\Delta^1$, cela signifierait que si dans $X$
on a une sous-catégorie fermée $X_1$ telle que
$X_1\lhook\joinrel\xrightarrow{i}X$ soit $W$-asphérique, alors
\deleted{\deleted{non seulement $i\in W$, mais aussi} $j:X_0\hookrightarrow X$ (imm. complémentaire) devrait être $\in W$. Prenons par contre
$X=X_0\cup\{e\}$, $e$ un ``objet initial
rajouté à $X_0$'', donc $X_1=\{e\}$. Alors \deleted{\unsure{gt} $X_1\sim X_1$,
et $X_1\mathrel{\unsure{\text{---}|\longrightarrow}}X$ est $W$-asph. sss $X_1$ est $W$-asph.}
\deleted{Et} $X_1=\{e\}\longrightarrow X$ est bien asph., \deleted{\unsure{mais}} $X$ est
$W$-asphérique, et $X_0\hookrightarrow X$ n'est dans $W$
\uline{que} si $X_0$ est elle-$=$ $W$-asphérique. Donc on
a notre contreexemple. Mais que dire de
$\underline{\mathrm{Prop}_0\,S}$, si \unsure{,,} $X$ est un $W$-\uline{fibré} \uline{propre} \inserted{sur $S$}, $\Phi(X)$ est-il encore propre ? \deleted{Le cas crucial devrait être \unsure{encore} $S=\Delta^1$, \unsure{or}}}
${}_{1\backslash}X\longrightarrow{}_{0\backslash}X$ est dans $W$ -- or
${}_{0\backslash}X\simeq X$, ${}_{1\backslash}X\simeq X_1$, la flèche est $X_1\mathrel{\unsure{\text{---}|\longrightarrow}}X$, qui est
bien dans $W$ étant $W$-asphérique ! C'est étrange et
\unsure{quasiment|carrément} \unsure{immoral|innocent}.* 
% slanted
\margin{On observe \unsure{plein les yeux}, c'est ainsi chaque fois que $S$ a un objet final}
Mais à force de gribouiller
je viens de me convaincre que même si
$X$ est cofibrée sur $S$ (donc aussi \inserted{$W$-}propre), $\Phi(X)$
n'est pas \inserted{$W$-}propre (ce qui impliquerait $W$-parfaite), ni
\scanpage{242}
\mspage{188}
% archive sheet number: 241
même un $W$-fibré, ce qui signifierait que
\[
{}_{s'\backslash}X\longrightarrow{}_{s\backslash}X
\qquad\text{associé à}\qquad s\mathrel{\unsure{\text{---}|\longrightarrow}}s'\quad\text{dans }S
\]
serait toujours dans $W$. Posant $Z'={}_{s'\backslash}S$, $Z={}_{s\backslash}S$,
cela signifierait que si on a un \inserted{co-}\uline{\unsure{isom|isom.} local}
\[
 i:Z'\longrightarrow Z
\]
où $Z',Z$ ont \deleted{des} \inserted{$\unsure{z'_0,z_0|z'_0\;z_0}$} éléments initiaux, et si
$X$ est cofibrée sur $Z$, $X'$ son image inverse sur
$Z'$, alors $X'\mathrel{\unsure{\text{---}|\longrightarrow}}X$ est dans $W$. \deleted{\deleted{\unsure{Voici un}
contre-exemple, soit} $Z=\Delta^2=\{0,1,2\}$, $Z'=\Delta^1=\{0,1\}$,
$i$ l'inclusion évidente (identique), qui est une
immersion fermée donc un coisom. local. La
donnée de $X$ cofibrée et \unsure{coscindée} sur $\deleted{S}\inserted{Z}=\Delta^2$
équivaut à une suite de deux foncteurs
\[
X_0\mathrel{\unsure{\text{---}|\longrightarrow}}X_1\mathrel{\unsure{\text{---}|\longrightarrow}}X_2
\]
(où $X_i=$ fibre de $X$ en $i$), $X'$ \unsure{estant|étant} \unsure{décrit|décrite}
alors par $X_0\longrightarrow X_1$. Comme $X$ $W$-propre \uline{,,} sur
$S$, a un élément final $e=2$, l'inclusion $X_e\mathrel{\unsure{\text{---}|\longrightarrow}}X_{/e}=X$
est $\in W$, donc dans $\mathrm{Hot}_W$ \deleted{$X$ s'identifie à} $X_e=X_2$
\deleted{de $=$ $X'$ s'identifie à $X_1$}} \unsure{toute fois|toutefois}, si $Z$ a \inserted{aussi} un
élément \uline{final} $e$, (comme c'était le cas de départ
pour $\Delta^1$) alors $\deleted{\unsure{X\simeq X_e}}$ $Z'\simeq{}_{z_0\backslash}Z$ en a aussi
un $e'$, qui s'envoie sur $e$, et \deleted{\unsure{de même}} on aura
le diagramme commutatif
\diagram{\begin{tikzcd}[ampersand replacement=\&,column sep=small,row sep=small]
X'\arrow[r,"j"]\&X\\
X'_{e'}\arrow[u,"W"]\arrow[r,"\sim"]\&X_e\arrow[u,"W"']
\end{tikzcd}}
, d'où $j\in W$. Donc pour faire un contreexemple,
